\chapter{Architecture of the Formalization}\label{ch:architecture}

\Cref{ch:leanprimer} taught the student to read a single Lean statement. This chapter
answers the next question: \emph{where do the statements quoted in this book live, how are
they organized, how much of the physics do they really capture, and what has to happen before
a theorem is allowed to carry the badge \TA?} The answer is an inventory and a procedure. The
inventory is a single Lake project containing seven libraries, five of them written without
Mathlib and two on top of it, whose compiled theorems number $326+99=425$ at the audit of
2026-09-17. The procedure is a chain of gates---build, \lean{sorry} count, axiom audit,
statement lock, and the rule that whoever produced a proof never grades it---each of which
exists because of a specific failure the project recorded. A student should care for two
reasons. First, the honest reading of any theorem in this book depends on knowing which
library it comes from: the older libraries model physical quantities by integers and prove
arithmetic consequences, while the newest one proves statements about real matrices of every
size. Second, the gates are the answer to the question every physicist asks of a proof
assistant---``if the machine checked it, what exactly can still go wrong?''---and the answer
is more interesting than ``nothing''.

%=====================================================================================
\section{One project, seven libraries}\label{sec:architecture:libraries}

A Lean 4 project is described by a file \lean{lakefile.lean}\index{lakefile}, read by the
build tool Lake. The project's file declares one package, pins one external dependency, and
registers seven libraries \source{lakefile.lean, ll.~1--53}. Three facts from that file
matter for everything that follows.

\begin{enumerate}[leftmargin=*,itemsep=2pt]
\item \textbf{The toolchain and Mathlib are pinned to the same tag.} The project uses Lean
  \lean{v4.33.1} and requires Mathlib at the git tag \lean{v4.33.1}, whose own
  \lean{lean-toolchain} file names exactly that compiler \source{lakefile.lean,
  ll.~12--17}. The comment in the file records why this is not pedantry: an earlier
  document had pinned a commit belonging to the \lean{v4.33.0} tag, which would have
  produced a cache of compiled Mathlib files unusable by the project's compiler. A
  formalization is reproducible only if the reader can rebuild it, and rebuilding Mathlib
  from source takes hours; the pin is what makes the pre-built cache usable.
\item \textbf{Two compiler limits are raised.} The package sets \lean{maxHeartbeats} to
  $10^6$ (the default is $2\times10^5$) and \lean{maxRecDepth} to $8000$
  \source{lakefile.lean, ll.~6--10}. Heartbeats are Lean's unit of elaboration work; a proof
  that exceeds the budget is rejected with a timeout, not accepted. \Cref{sec:architecture:gates}
  returns to a consequence of this setting that cost the project several hours.
\item \textbf{Five libraries are default targets; two are not.} A plain \lean{lake build}
  compiles the five libraries in the upper rows of \cref{tab:architecture:libraries}, none
  of which imports Mathlib \source{lakefile.lean, ll.~19--39}. The Mathlib-based libraries
  \lean{StringTheoryFormalization} and \lean{DualScaleStream2} must be named explicitly,
  as in \lean{lake build} \lean{DualScaleStream2}. This is by design: the file says the
  second was kept off the default list ``so a Stream 2 failure never turns the green
  Stream 1 core red'' \source{lakefile.lean, ll.~41--53}.
\end{enumerate}

\Cref{tab:architecture:libraries} lists the seven libraries with the numbers that the
project's own audit reports. The theorem counts are those of \lean{tools/axiom\_audit.py}
run on 2026-09-17, quoted from the revision brief \source{PAPER\_REVISION\_BRIEF §1,
ll.~25--35}; the module counts are the number of \lean{.lean} files imported by each
library's root file, counted by the author of this chapter on the same checkout.

\begin{table}[ht]\centering\small
\caption{The seven libraries. ``Mathlib'' says whether any module imports Mathlib.
``Imports'' lists the other project libraries the library depends on. Theorem counts are
the audited \lean{theorem}/\lean{lemma} declarations \source{PAPER\_REVISION\_BRIEF §1,
ll.~25--35}.}\label{tab:architecture:libraries}
\begin{tabular}{@{}lccc>{\raggedright\arraybackslash}p{4.6cm}@{}}\toprule
Library & Modules & Theorems & Mathlib & Imports (project libraries) \\ \midrule
\lean{StringTheoryFoundation}    & 16 & 56 & no  & --- \\
\lean{DualScaleM24Formalization} & 11 & 61 & no  & --- \\
\lean{DoubleFieldTheory}         & 7  & 44 & no  & --- \\
\lean{DualScaleValidation}       & 4  & 23 & no  & \lean{DoubleFieldTheory}, \lean{DualScaleM24Formalization} \\
\lean{Lean5Corpus}               & 12 & 53 & no  & the three above and \lean{StringTheoryFoundation} \\
\lean{StringTheoryFormalization} & 35 & 89 & yes & \lean{StringTheoryFoundation} \\
\lean{DualScaleStream2}          & 19 & 99 & yes & \lean{StringTheoryFormalization} \\ \midrule
Stream 1 (first six)             & 85 & 326 & & 3353 build jobs \\
Stream 2 (last)                  & 19 & 99  & & 3670 build jobs \\
\bottomrule\end{tabular}
\end{table}

The words \emph{Stream 1} and \emph{Stream 2}\index{Stream 1, Stream 2} are the project's
names for the two generations of the corpus: Stream 1 is everything that existed before
2026-09-16, re-gated in the same session; Stream 2 is the library \lean{DualScaleStream2} written
during the tiered proving run of 2026-09-16/17 described in \cref{ch:pipeline}
\source{STREAM2\_WORKFLOW.md §1, ll.~7--18}. The job counts are what Lake reports when it
compiles a target together with everything it depends on; the 3670 jobs of Stream 2 include
the parts of Mathlib it imports, which is why a nineteen-module library needs more jobs than
the eighty-five modules of Stream 1.

\begin{figure}[ht]\centering
\begin{tikzpicture}[
  lib/.style={draw,rounded corners=2pt,minimum height=7mm,inner sep=4pt,font=\small\ttfamily,align=center},
  core/.style={lib,fill=green!8,draw=tierA},
  ml/.style={lib,fill=blue!6,draw=blue!45!black},
  ext/.style={lib,fill=black!6,draw=black!55,font=\small\sffamily},
  imp/.style={-{Stealth[length=2mm]},thick}]
  \node[core] (STF)  at (0,4.2)   {StringTheoryFoundation\\[-2pt]{\scriptsize\sffamily 16 modules, 56 thm}};
  \node[core] (DSM)  at (5.2,4.2) {DualScaleM24Formalization\\[-2pt]{\scriptsize\sffamily 11, 61}};
  \node[core] (DFT)  at (10.2,4.2){DoubleFieldTheory\\[-2pt]{\scriptsize\sffamily 7, 44}};
  \node[core] (DSV)  at (7.7,2.4) {DualScaleValidation\\[-2pt]{\scriptsize\sffamily 4, 23}};
  \node[core] (L5C)  at (3.6,2.4) {Lean5Corpus\\[-2pt]{\scriptsize\sffamily 12, 53}};
  \node[ml]   (STFo) at (0,0.6)   {StringTheoryFormalization\\[-2pt]{\scriptsize\sffamily 35, 89}};
  \node[ml]   (DSS2) at (5.2,0.6) {DualScaleStream2\\[-2pt]{\scriptsize\sffamily 19, 99}};
  \node[ext]  (MLB)  at (10.2,0.6){Mathlib v4.33.1};
  \draw[imp] (DSV) -- (DFT);  \draw[imp] (DSV) -- (DSM);
  \draw[imp] (L5C) -- (DSV);  \draw[imp] (L5C) -- (STF); \draw[imp] (L5C) to[bend right=12] (DFT);
  \draw[imp] (STFo) -- (STF); \draw[imp] (STFo) -- (MLB);
  \draw[imp] (DSS2) -- (STFo); \draw[imp] (DSS2) -- (MLB);
  \draw[dashed,black!60] (-2.2,1.55) -- (12.2,1.55);
  \node[font=\scriptsize\sffamily,anchor=west,black!70] at (-2.2,1.85) {Mathlib-free core (default targets)};
  \node[font=\scriptsize\sffamily,anchor=west,black!70] at (-2.2,1.25) {Mathlib-based (built by name)};
\end{tikzpicture}
\caption{Import structure of the seven libraries. An arrow $A\to B$ means some module of
$A$ imports a module of $B$. Above the dashed line nothing imports Mathlib; the two
libraries below it do. \lean{DualScaleStream2} imports five modules of
\lean{StringTheoryFormalization} and nothing else in the project; no library imports
\lean{DualScaleStream2}.}\label{fig:architecture:imports}
\end{figure}

\Cref{fig:architecture:imports} shows the import graph, obtained by grepping every
\lean{import} line in the checkout. Two features deserve comment. The three libraries in the
top row import nothing but themselves: each is self-contained, which is why they could be
written before Mathlib was available to the project at all. The Mathlib-based libraries form
a chain: \lean{DualScaleStream2} imports exactly five modules of
\lean{StringTheoryFormalization} (the circle mass spectrum, the Hodge numbers, the Narain
lattice, the K3 signature theorem and the $M_{24}$ table), which in turn imports three
modules of \lean{StringTheoryFoundation}. The workflow document states the
design rule behind the chain: ``Stream 2 lives in \lean{DualScaleStream2}, imports Stream 1,
never the reverse; Stream 1's green core cannot be turned red by Stream 2 work''
\source{STREAM2\_WORKFLOW.md §2, ll.~32--33}. The rule is enforced by the import direction
itself: a broken file in \lean{DualScaleStream2} cannot stop any other library from
compiling, because nothing depends on it.

\begin{physicsbox}[Why cross-library imports are the interesting ones]
An import from \lean{DualScaleStream2} into Stream 1 is not bookkeeping; it is a theorem
that two independently written pieces of mathematics agree. The module
\lean{DFT/GeneralizedMetric.lean} imports Stream 1's \lean{TDualityMassSpectrum} in order to
prove that the doubled mass form $Z^{\mathsf T}\HH Z$ on a circle equals half of Stream 1's
$P_L^2+P_R^2$ (\cref{ch:genmetric}); \lean{Lattice/Hyperbolic.lean} imports
\lean{NarainLattice} to prove that Stream 2's hyperbolic plane $U$ is \emph{equal} to Stream
1's Narain Gram matrix; \lean{Lattice/K3T2Signature.lean} imports \lean{K3SignatureTheorem}
to prove that the lattice signature $(3,19)$ agrees with the Hodge-theoretic
$(2h^{2,0}+1,\,h^{1,1}-1)$ \source{PAPER\_REVISION\_BRIEF §3.1, ll.~78--84}. Each such
theorem is a consistency check between two formalizations of the same physics, checked by
the kernel. A physicist would call them cross-calibrations.
\end{physicsbox}

%=====================================================================================
\section{What each library contains}\label{sec:architecture:contents}

The following inventory names each library's modules and says, in one or two sentences,
what kind of mathematics each contains. It is deliberately flat: the depth of the
statements, which is what determines how a theorem may be quoted, is the subject of
\cref{sec:architecture:depth}.

\paragraph{\lean{StringTheoryFoundation} (16 modules, 56 theorems).} The oldest library.
Its modules are grouped by physical topic---\lean{Core/Topology} (Euler characteristics of
$K3$ and $T^2$ from Betti-number records), \lean{Duality/T\_Duality} (the exchange
$(n,w)\mapsto(w,n)$ of integer quantum numbers as an involution), \lean{Duality/DualScale}
(dual length pairs as integer relations), \lean{K3/K3Surfaces} (Betti numbers),
\lean{StringTheory/K3xT2}, \lean{StromingerSYZ}, \lean{Swampland}, \lean{VafaSwampland},
\lean{TadpoleCancellation}, \lean{WittenDuality}---together with four modules named after
external formalizations (Navier--Stokes, Fermat and modular forms, statistical learning,
ATLAS geometry; listed and described in \cref{sec:architecture:bridges}) and two
physics-library bridges (\lean{PhysLibKinematicsBridge}, \lean{TensorNetworkBridge}).

\paragraph{\lean{DualScaleM24Formalization} (11 modules, 61 theorems).} The library of the
programme's own claims: \lean{DualScale/EffectiveMetric} (the effective scale
$R+\ap/R$ and its lower bound, \cref{ch:dualscale}), \lean{DualScale/Sym2Lock},
\lean{Moonshine/MathieuRigidity} (the ratio $77/60$ and the identity $462\cdot60=360\cdot77=27720$,
\cref{ch:moonshine}), \lean{Moonshine/KummerTadpole}, \lean{Moonshine/GysinBEMSequence}, the
four \lean{FrontierTriad} modules (\cref{ch:triad}), and the two \lean{Epistemic} modules,
\lean{Tier} and \lean{Ledger}, which formalize the tier system itself and are the subject
of \cref{sec:architecture:epistemic}.

\paragraph{\lean{DoubleFieldTheory} (7 modules, 44 theorems).} Integer-valued
$2\times2$ matrix models of the DFT objects of \cref{ch:doubled,ch:genmetric,ch:courant,ch:dftaction}:
\lean{GeneralizedGeometry} (a structure \lean{Mat2} with four integer entries, the pairing
$\eta$, a $B$-twist, a generalized metric), \lean{CourantAlgebroid}, \lean{ActionCurvature},
\lean{TDualityBuscher}, \lean{K3Topology}, \lean{TorusMoonshine} and a \lean{PhysicsDSL}.

\paragraph{\lean{DualScaleValidation} (4 modules, 23 theorems).} Three ``use case''
modules and an \lean{Observables} module of exact rational identities and window
inequalities. The module header of \lean{Observables} is unusually explicit about its own
scope: ``Tier A for its arithmetic only \ldots\ not a derivation of any physical observable
from first principles'', with the numerical relations labelled Tier C
\source{Observables.lean, ll.~9--24}. \Cref{sec:architecture:epistemic} returns to the
lesson in that header.

\paragraph{\lean{Lean5Corpus} (12 modules, 53 theorems).} A \lean{Foundations/Narrative}
module and eleven \lean{Problems/Problem$k$} modules, one per named problem (Navier--Stokes
helicity, Mathieu--Frobenius rigidity, dual-scale TCC, Mukai monodromy, Kolmogorov cascade,
flux swampland, Courant torsion, Golay holography, Kummer modularity, SYM instanton,
entanglement entropy). It imports the three libraries above and \lean{StringTheoryFoundation}
and is the only library that ties the Mathlib-free core together.

\paragraph{\lean{StringTheoryFormalization} (35 modules, 89 theorems).} The first
Mathlib-based library. Its \lean{Frontier} modules (central charge, chiral primaries,
Hodge numbers, moduli geodesics, $SL(2,\C)$ symmetry, F-term potential), its
\lean{StringDynamics} modules (Kummer blow-up, tadpole constraint, the $M_{24}$ table,
BPS multiplicities, Mukai lattice, T-duality Gysin, vertex operators, and several numerical
or search-related modules), its five \lean{UseCases} (critical dimension, K3 signature,
Mathieu tower, Narain lattice, T-duality mass spectrum) and its \lean{NSMath} modules
(fractional Sobolev spaces, Fourier multipliers, mild PDEs, energy bounds). Several of its
theorems are quoted in earlier chapters, for instance the circle mass spectrum of
\cref{ch:tduality} and the Narain Gram matrix of \cref{ch:narain}. Its header history is
instructive: the library ``had 8 real build failures despite being self-documented as
26--28/30 modules, complete'' when the checkout was rebuilt on 2026-09-16
\source{LL.md §0, ll.~148--155}.

\paragraph{\lean{DualScaleStream2} (19 modules, 99 theorems).} The library written by the
tiered pipeline, in six groups that mirror the workflow's phases \source{STREAM2\_WORKFLOW.md
§5, ll.~67--76}: \lean{Lattice} (\lean{Basic}, \lean{E8}, \lean{E8PosDef}, \lean{Hyperbolic},
\lean{K3T2Signature}, \lean{Mukai}, \lean{Reflection}; \cref{ch:lattices,ch:k3lattice,ch:mukai}),
\lean{TDuality} (\lean{ODD}, \lean{Factorized}, \lean{Spectrum}, \lean{Mirror},
\lean{SL2Product}; \cref{ch:odd,ch:torus2}), \lean{DFT} (\lean{GeneralizedMetric},
\lean{BShift}, \lean{SectionCondition}; \cref{ch:genmetric,ch:doubled}),
\lean{DualScale/TraceBound} (\cref{ch:dualscale}), \lean{Flux} (\lean{Tadpole},
\lean{Integrality}; \cref{ch:tadpole}) and \lean{Moonshine/EOT} (\cref{ch:moonshine}). Every
module header states the physical background, the pinned literature sources with line
numbers, what is not proved, and the proof techniques used.

%=====================================================================================
\section{Depth: arithmetic shadows and faithful linear algebra}\label{sec:architecture:depth}

Two libraries in the project both contain a theorem that ``the generalized metric satisfies
$\HH\eta\HH=\eta$''. Read side by side, they show exactly what the phrase \emph{arithmetic
shadow}\index{arithmetic shadow} means in this book, and why the library a theorem comes
from changes how it may be quoted.

In \lean{DoubleFieldTheory/GeneralizedGeometry.lean} the generalized metric is a
$2\times2$ matrix with integer entries, built from an integer $g$ and an integer standing for
$g^{-1}$ \source{DoubleFieldTheory/GeneralizedGeometry.lean, ll.~189--209}:
\begin{leancode}
def GenMetric (g : Int) (ginv : Int) : Mat2 :=
  { a := g, b := 0, c := 0, d := ginv }
theorem gen_metric_symmetric (g ginv : Int) :
    (GenMetric g ginv).b = (GenMetric g ginv).c := by rfl
theorem gen_metric_duality (g : Int) (h : g = 1) :
    MatMul (GenMetric g g) (MatMul ODD_Eta (GenMetric g g)) = ODD_Eta := by ...
\end{leancode}
Read literally: \lean{GenMetric g ginv} is the diagonal matrix $\diag(g,g^{-1})$ with the
inverse supplied by hand, and the ``duality'' theorem says that $\HH\eta\HH=\eta$ holds for
$\HH=\diag(1,1)$, under the hypothesis $g=1$. Both proofs are \lean{rfl}: the statements
are true by computation on literals. This is the $d=1$, $B=0$, self-dual-radius corner of
the physics, and only that corner. It is a correct theorem, and it is Tier A, but it is a
theorem about the integer $1$.

In \lean{DualScaleStream2/DFT/GeneralizedMetric.lean} the same object is a block matrix
over $\R$ of size $2d\times2d$ for every $d$, built from an arbitrary $G$ and $B$
\source{DFT/GeneralizedMetric.lean, ll.~91--92, 103--104}:
\begin{leancode}
noncomputable def genMetric (G B : Matrix (Fin d) (Fin d) ℝ) :
    Matrix (Charge d) (Charge d) ℝ :=
  fromBlocks (G - B * G⁻¹ * B) (B * G⁻¹) (-(G⁻¹ * B)) G⁻¹
theorem etaR_genMetric_sq (G B : Matrix (Fin d) (Fin d) ℝ) (hG : IsUnit G.det) :
    (etaR d * genMetric G B) * (etaR d * genMetric G B) = 1 := by ...
\end{leancode}
Here $G^{-1}$ is Mathlib's matrix inverse, the hypothesis is that $G$ is invertible, $B$ is
arbitrary, and the conclusion $(\eta\HH)^2=1$ is equivalent to $\HH\eta\HH=\eta$ because
$\eta^2=1$. This is the statement a physicist means (\cref{ch:genmetric}): it covers every
torus dimension, every metric and every $B$-field. The proof is genuine block-matrix
algebra with inverses, and it needed the most capable model tier in the pipeline
\source{STREAM2\_WORKFLOW.md §6, ll.~136--143}.

\begin{definition}[Arithmetic shadow]\label{def:architecture:shadow}
A Lean theorem is an \emph{arithmetic shadow} of a physical statement when its objects are
integers, rationals or fixed small matrices standing in for the physical quantities, so that
the theorem is a numerical instance or a consistency check of the physical statement rather
than the statement itself. Five of the seven libraries---\lean{DoubleFieldTheory},
\lean{StringTheoryFoundation}, \lean{DualScaleM24Formalization}, \lean{Lean5Corpus},
\lean{DualScaleValidation}---consist mostly of shadows in this sense; this is how the book's
bible instructs authors to describe them.
\end{definition}

A shadow is not worthless. The identity $\chi(K3)=1-0+22-0+1=24$ is a shadow of the
Betti-number computation, and it is exactly the arithmetic a student does on paper. The
Burnside check $\sum_i \dim(\rho_i)^2 = |M_{24}| = 244\,823\,040$ over the 26 irreducible
dimensions is a shadow, and it caught a wrong table that had been in the repository for
months (\cref{sec:architecture:gates}). What a shadow cannot do is stand in for the general
theorem. To turn \lean{gen\_metric\_duality} into \lean{etaR\_genMetric\_sq} one needs, in
order: a type of $d\times d$ real matrices with a genuine inverse (in Mathlib,
\lean{Matrix.nonsing\_inv\_mul}: $A^{-1}A=1$ when $\det A$ is a unit), block matrices (\lean{Matrix.fromBlocks}), the invertibility
hypothesis in place of $g=1$, and a proof that is a computation with the algebraic laws of
inverses rather than a computation on literals. That is precisely the gap Mathlib fills, and
it is why the two Mathlib-based libraries are where the faithful statements live.

\begin{pitfallbox}[A theorem about a fixed matrix is not a theorem about all lattices]
The temptation runs the other way too. \lean{DualScaleStream2/Lattice/E8.lean} proves that
one explicit $8\times8$ integer matrix is symmetric, even and unimodular; it does not prove
that ``the $E_8$ lattice'' has these properties, because no lattice, no basis change and no
isomorphism appears in the file (\cref{ch:lattices}). Similarly the signature arithmetic in
\lean{K3T2Signature} adds up signatures of named summands; that $H^2(K3,\Z)$ \emph{is}
$E_8(-1)^{\oplus2}\oplus U^{\oplus3}$ is Tier L, pinned to Huybrechts
\source{PAPER\_REVISION\_BRIEF §3.1, ll.~81--87}. The rule of thumb: the tier badge attaches
to the Lean statement as written, and the physical name in the docstring is a promise the
kernel never checked.
\end{pitfallbox}

\begin{example}[Reading the depth of a statement]\label{ex:architecture:depth}
A student wants to cite ``the $E_8$ Cartan matrix is positive definite'' with a Tier A
badge. She finds three candidate statements. (i) \lean{cartanE8\_det} in
\lean{DualScaleStream2/Lattice/E8PosDef.lean}, which says that the integer determinant of
the explicit matrix is $1$. (ii) \lean{cartanE8\_posDef} in the same file, whose statement is
Mathlib's \lean{Matrix.PosDef} predicate applied to the explicit matrix over $\R$, proved
through the sum-of-squares identity $x^{\mathsf T}E_8x=\sum_k D_k y_k^2$ with pivots
$2,\tfrac32,\tfrac43,\ldots,\tfrac18$ \source{PAPER\_REVISION\_BRIEF §3.1, ll.~75--77}.
(iii) A sympy script that computes the eigenvalues. Only (ii) is the Tier A statement she
wants; (i) is a shadow of $\det E_8=1$, which does not imply positive definiteness; (iii) is
a symbolic check and by the book's rules is never Tier A. The product of the pivots is
$2\cdot\tfrac32\cdot\tfrac43\cdot\tfrac54\cdot\tfrac65\cdot\tfrac76\cdot\tfrac87\cdot\tfrac18
= 8\cdot\tfrac18 = 1$, which is the determinant---a two-line check the student can do herself
before trusting the file.
\end{example}

%=====================================================================================
\section{The peripheral bridge modules}\label{sec:architecture:bridges}

Four modules of \lean{StringTheoryFoundation} carry the names of external formalization
projects:
\begin{itemize}[leftmargin=*,itemsep=0pt]
\item \lean{FluidDynamics/NavierStokesBridge} (a Navier--Stokes formalization),
\item \lean{ModularForms/FermatModularBridge} (a Fermat--modularity formalization),
\item \lean{StatisticalLearning/StatisticalLearningBridge} (a statistical-learning library),
\item \lean{Atlas/AtlasGeometryBridge} (an autoformalized textbook library).
\end{itemize}
The projects they name are vendored, read-only, as git submodules under
\lean{lean4basesource/}. This section states what the four modules contain, because their
names promise more than their code delivers and the project's own review says so in each
file.

Every one of the four carries a ``scope note'' inserted after review. The note in the
Navier--Stokes module is representative: the file ``does \textbf{not} import anything from
that project---check the \lean{import} lines above; there is only
\lean{StringTheoryFoundation.Core.Topology}, this project's own file. The theorems below are
self-contained \lean{Nat}/\lean{Int} arithmetic named after, and inspired by, the cited
external work, not a machine-checked bridge to it'' \source{NavierStokesBridge.lean,
ll.~40--47}. The Fermat and statistical-learning modules carry the same sentence
\source{FermatModularBridge.lean, ll.~42--49; StatisticalLearningBridge.lean, ll.~40--44},
and the Atlas module notes that it ``has \textbf{no} \lean{import} statement at all''
\source{AtlasGeometryBridge.lean, ll.~15--21}. The workflow document records the technical
reason no genuine bridge is possible at present: the external corpora are pinned to other
Lean versions (\lean{atlas-lean} at v4.29.0, the Navier--Stokes project at v4.34.0-rc2) and
``cannot be \lean{require}d at this project's v4.33.1 pin'' \source{STREAM2\_WORKFLOW.md §2,
ll.~34--37}. \Cref{tab:architecture:bridges} lists what each module actually proves.

\begin{table}[ht]\centering\small
\caption{The four bridge modules and their theorems, read literally from the source. All are
Tier A as arithmetic; none imports the external project it is named after.}
\label{tab:architecture:bridges}
\begin{tabular}{@{}p{3.8cm}>{\raggedright\arraybackslash}p{5.5cm}>{\raggedright\arraybackslash}p{5.2cm}@{}}\toprule
Module & Theorems & What is literally proved \\ \midrule
\lean{NavierStokesBridge} \newline\source{ll.~53--111} &
\lean{zero\_mode\_laplacian}, \lean{fundamental\_mode\_positive},
\lean{energy\_dissipation\_monotonic}, \lean{sound\_speed\_is\_luminal} &
$0^2+0^2=0$ and $1^2+0^2=1$ for an integer wave-vector; $e-d\le e$ for natural numbers
$e,d$ (truncated subtraction); $1/(2-1)=1$ in $\N$. \\[3pt]
\lean{FermatModularBridge} \newline\source{ll.~60--113} &
\lean{kummer\_fixed\_points\_dim4}, \lean{exceptional\_divisor\_cartan\_a1},
\lean{mukai\_lattice\_rank\_equals\_24}, \lean{mukai\_signature\_difference} &
$2^4=16$; the constant $-2$ equals $-2$; $1+22+1=24$; $4-20=-16$. \\[3pt]
\lean{StatisticalLearningBridge} \newline\source{ll.~46--81} &
\lean{zero\_empirical\_risk\_}\allowbreak\lean{on\_perfect\_proofs},
\lean{rademacher\_bound\_positive} &
$(n-n)\cdot10^4/n=0$ for $n>0$; $2k>0$ for $k>0$. \\[3pt]
\lean{AtlasGeometryBridge} \newline\source{ll.~23--78} &
\lean{atlas\_k3\_euler\_characteristic}, \lean{atlas\_k3\_signature},
\lean{atlas\_poincare\_curvature\_negative} &
$1-0+22-0+1=24$; $3-19=-16$; the constant $-1$ equals $-1$. \\
\bottomrule\end{tabular}
\end{table}

Three of these facts are genuine, if elementary, shadows of physics the book has met:
$\chi(K3)=24$ (\cref{ch:k3}), the signature $-16$ of the intersection form and $(4,20)$ of
the Mukai lattice (\cref{ch:k3lattice,ch:mukai}), and the $16$ fixed points of the Kummer
involution (\cref{ch:kummer}). The others are definitional restatements, and the files say
so: the comment above \lean{exceptional\_divisor\_cartan\_a1} reads ``This is a definitional
restatement, not a derivation; recorded for cross-reference from other modules rather than
as an independent check'' \source{FermatModularBridge.lean, ll.~72--76}, and a theorem
that had checked $24=24$ was removed \source{FermatModularBridge.lean, ll.~99--103}. The
Navier--Stokes energy theorem once took $e_2\le e_1$ as a hypothesis and returned it as the
conclusion; it now derives $e-d\le e$ from \lean{Nat.sub\_le}
\source{NavierStokesBridge.lean, ll.~87--97}.

The honest summary is the one the project's own review reached. These modules are
\emph{documents of intent}: each header describes a real piece of mathematics (the
fluid--gravity correspondence, the modularity theorem, PAC generalization bounds,
autoformalized four-manifold topology) and a real external formalization, and each body
proves a handful of integer identities that would survive in any faithful version. They are
the smallest shadows in the corpus, they are counted among the 56 audited theorems of
\lean{StringTheoryFoundation}, and no chapter of this book cites them as evidence for the
physics in their headers. A faithful bridge would begin by importing the external
declarations that the Navier--Stokes file's own ``verified pointer'' names
\source{NavierStokesBridge.lean, ll.~48--51}, which the pinned toolchains currently forbid.

%=====================================================================================
\section{The gates}\label{sec:architecture:gates}

A theorem enters this book with the badge \TA{} only after passing a fixed sequence of
checks. The workflow document states the accept criterion in one sentence
\source{STREAM2\_WORKFLOW.md §2, ll.~22--24}:
\begin{quote}
``A goal is closed when \lean{lake build} \lean{DualScaleStream2} is green, the
zero-\lean{sorry} audit passes, and \lean{\#print axioms} shows only\\
\lean{propext}, \lean{Classical.choice}, \lean{Quot.sound}. No \lean{native\_decide}.''
\end{quote}
The per-package pipeline names the gates G1--G4 \source{STREAM2\_WORKFLOW.md §4,
ll.~52--65}; \cref{tab:architecture:gates} collects them with the failure each one was
built to catch. We take them in order.

\begin{table}[ht]\centering\small
\caption{The gates, the tool that implements each, and the recorded failure that motivated
it. Line references are to \lean{STREAM2\_WORKFLOW.md} (W) and \lean{LL.md} (L).}
\label{tab:architecture:gates}
\begin{tabular}{@{}p{2.6cm}p{4.2cm}p{7.6cm}@{}}\toprule
Gate & Tool & Catches \ldots \\ \midrule
Build (G1) & \lean{lake build <lib>} & files that do not elaborate; a self-reported
``26--28/30 modules'' that was 8 real failures \source{L §0.2, ll.~168--177} \\[2pt]
\lean{sorry} count & regex after stripping comments & goals left open; a report claiming ``5
clean \lean{sorry}s'' that had 8 compile errors \source{W §6, ll.~153--157} \\[2pt]
Statement review (G2) & a human or orchestrator reads each statement & mis-formalization,
``the one step no kernel can check'' \source{W §4, l.~59} \\[2pt]
Axiom audit (G3) & \lean{tools/axiom\_audit.py} & \lean{sorryAx} inherited through
dependencies; \lean{native\_decide} in a ``clean'' library \source{L §S2.5, ll.~56--61} \\[2pt]
Statement lock & \lean{tools/statement\_lock.py} & a statement weakened between review and
proof \source{L §S2.10, ll.~87--92} \\[2pt]
Producer $\ne$ verifier & the gate runs separately from the prover & agents misstating their
own compile status \source{L §S2.3, ll.~39--45} \\[2pt]
Record (G4) & declaration index, graph, ledger entry pinned to a commit hash & counts that
drift from the code \source{W §4, l.~64} \\
\bottomrule\end{tabular}
\end{table}

\subsection{Build and \lean{sorry} count}\label{sec:architecture:build}

The build is the gate everyone expects, and the project's log shows why it must be re-run
rather than believed. Lesson 0.2 of the log records that a library documented as complete
had eight real failures once ``the full dependency chain was exercised'', and draws the
general rule: ``after any claim of `X builds clean' from a document, memory note, or your own
earlier turn \ldots\ re-run the actual build before trusting it'' \source{LL.md §0.2,
ll.~168--177}. Lesson 0.3 adds a subtlety a student of any compiled language will
recognize: fixing one file can \emph{unblock} elaboration of files that depend on it and so
reveal failures that were always present but never reached, which is why the whole target
must be rebuilt after every fix \source{LL.md §0.3, ll.~179--186}.

The keyword \lean{sorry}\index{sorry@\lean{sorry}} closes any goal without a proof and makes
Lean emit a warning; a library can build green with hundreds of them. The count is taken by
a regular expression after stripping comments, so that a docstring mentioning the word does
not count \source{tools/axiom\_audit.py, ll.~30--34}. A \lean{sorry} is not a defect during
development---the pipeline of \cref{ch:pipeline} deliberately begins with statements whose
proofs are \lean{sorry}---but the count must be zero at the gate, and the audit of the next
subsection makes the count redundant by catching \lean{sorry} at the axiom level.

\subsection{The axiom audit}\label{sec:architecture:axioms}

The decisive check is \lean{\#print axioms}. Every Lean theorem, when printed this way,
lists the axioms its proof term ultimately depends on. Three are standard and harmless
(\cref{ch:leanprimer}): \lean{propext}, \lean{Classical.choice}, \lean{Quot.sound}. Two
others signal trouble. A theorem proved with \lean{sorry} depends on the axiom
\lean{sorryAx}, and so does every theorem that uses it, however indirectly: the axiom
propagates through the dependency graph, which a \lean{sorry} grep cannot follow. A theorem
proved by \lean{native\_decide} depends on \lean{Lean.ofReduceBool}, an axiom asserting that
the compiled evaluator agrees with the kernel; the kernel did not check the computation, the
compiler did.

The script \lean{tools/axiom\_audit.py} automates this over a whole library
\source{tools/axiom\_audit.py, ll.~1--102}. Its logic is short enough to describe fully.
\begin{enumerate}[leftmargin=*,itemsep=1pt]
\item Collect every \lean{theorem}/\lean{lemma} name from the library's source files,
  tracking \lean{namespace}/\lean{end} blocks to build the fully qualified name, after
  stripping comments (ll.~30--49).
\item Write a probe file that imports the built library and contains one
  \lean{\#print axioms} line per name; run it with \lean{lake env lean} (ll.~62--67).
\item If the probe itself fails to elaborate---typically because the library was not
  rebuilt---exit with code 2 and the message \texttt{AUDIT ERROR} (ll.~70--77). This clause
  was added after an audit of a stale build reported ``89 theorems audited, 89 failing
  (MISSING)'', which ``looks like a proof result, but it was an infrastructure failure''
  \source{LL.md §S2.11, ll.~94--99}.
\item Parse each theorem's axiom list, subtract the standard set
  \lean{\{propext, Classical.choice, Quot.sound\}}, and fail the gate if anything remains
  (ll.~79--98).
\end{enumerate}
The audit is validated the way an experimentalist validates a detector: with a positive
control (a library known to be clean) and a negative control (a library with known
\lean{sorry}s) before its verdicts are trusted \source{LL.md §S2.5, ll.~59--61}.

\begin{example}[Running the audit by hand on the epistemic module]\label{ex:architecture:audit}
The author of this chapter ran a five-line probe against the compiled library on the
checkout of 2026-09-17:
\begin{leancode}
import DualScaleM24Formalization.Epistemic.Ledger
#print axioms SocrateAI.Epistemic.tier_trans
#print axioms SocrateAI.Epistemic.tier_le_of_depends
#print axioms SocrateAI.Epistemic.no_kernel_claim_rests_on_weaker
#print axioms SocrateAI.Epistemic.witnessSound_sound
#print axioms SocrateAI.Epistemic.witnessUnsound_not_sound
\end{leancode}
The output was \lean{[propext, Quot.sound]} for the first three and \lean{[propext]} for
the two witnesses. All five pass: the sets are subsets of the standard three, and none is
\lean{sorryAx} or \lean{Lean.ofReduceBool}. The student should notice two things. First,
\lean{Classical.choice} is absent: these proofs (by \lean{omega}, induction and
\lean{decide}) never invoke the axiom of choice, and the audit records that fact for free.
Second, the
count of axioms is not a measure of strength---the witnesses use fewer axioms than the
general theorem only because \lean{decide} evaluates a finite case by computation.
\end{example}

\subsection{What the audit found in a ``verified'' library}\label{sec:architecture:found}

The reason the project trusts this gate is that it found real defects in code that the
build and the \lean{sorry} grep had passed for months.

\begin{historybox}[Two defects caught by the audit and by a guard theorem]
Both are recorded in the workflow document and the revision brief. (1)
\lean{StringDynamics.bps\_ratio\_reduced}, the theorem that $77/60$ is in lowest terms,
``had relied on \lean{native\_decide} (trusts compiled code, not the kernel); it is now a
kernel proof'' by \lean{norm\_num} \source{PAPER\_REVISION\_BRIEF §2, ll.~44--45;
STREAM2\_WORKFLOW.md §6, ll.~127--128}. (2) The table \lean{StringDynamics.M24RepDim} of the
26 irreducible dimensions of $M_{24}$ ``was wrong: 10395 and 483 duplicated, second 990
and third 1035 missing; squares summed to 351,061,029 $\neq |M_{24}|$'', because its only
theorem checked a single entry. It was replaced by the table of Eguchi--Ooguri--Tachikawa
and guarded by \lean{M24RepDim\_sum\_sq}, whose statement is
$\sum_{i<26}\lean{M24RepDim}(i)^2 = 244\,823\,040$ \source{STREAM2\_WORKFLOW.md §6,
ll.~129--132; MathieuM24.lean, ll.~118--121}. The lesson the log draws is general: ``guard
every data table with a theorem that would break if the table were wrong''
\source{LL.md §S2.6, ll.~63--67}.
\end{historybox}

The second defect deserves a moment's arithmetic, because it shows what a guard theorem
buys. The order of $M_{24}$ is $2^{10}\cdot3^3\cdot5\cdot7\cdot11\cdot23=244\,823\,040$
(\cref{ch:m24}). Burnside's identity $\sum_i\dim(\rho_i)^2=|G|$ holds for every finite group
(Tier L, standard; not checked against a source in this book's library). The wrong table
summed to $351\,061\,029$, a discrepancy of $106\,237\,989$, so a single theorem comparing
the sum to the group order rejects it. The kernel evaluates the sum by unfolding the 26
cases (\lean{simp [Fin.sum\_univ\_succ, M24RepDim]}), which is a shadow computation in the
sense of \cref{def:architecture:shadow}---but a shadow attached to the right invariant is a
strong test, exactly as a checksum is.

\subsection{The statement lock}\label{sec:architecture:lock}

The kernel checks that a proof proves its statement. It cannot check that the statement is
still the one that was reviewed. In a workflow where proofs are searched for by automated
agents this is the central vulnerability: ``An agent (or a person) can make a goal
`provable' by weakening it'' \source{tools/statement\_lock.py, ll.~5--7}. The tool that
closes the gap hashes statements \source{tools/statement\_lock.py, ll.~1--137}:
\begin{itemize}[leftmargin=*,itemsep=1pt]
\item for a \lean{theorem} or \lean{lemma}, the text from the keyword up to the top-level
  \lean{:=}, tracked with a bracket-depth counter so that a \lean{:=} inside a binder does
  not end the statement early (ll.~69--83); replacing \lean{:= by sorry} with a proof
  therefore leaves the hash unchanged;
\item for a \lean{def}, \lean{abbrev} or \lean{structure}, the whole declaration including
  its body, ``changing a definition changes the meaning of every statement that uses it''
  (ll.~9--13);
\item after removing comments with a small scanner that handles Lean's nested block
  comments, and normalizing whitespace, so documentation can change freely (ll.~46--66).
\end{itemize}
The first 16 hex digits of the SHA-256 of the normalized text are stored per file and per
declaration in \lean{docs/statement\_lock.json}; \lean{--update} is run after statement
review, \lean{--check} before any proof is accepted, and a changed or vanished locked
declaration fails the check while a new one is only reported (ll.~17--22, 117--133). On the
checkout used for this chapter the lock covers 54 files and 355 declarations, all in the
two Mathlib-based libraries; the Mathlib-free core is not locked.

\begin{example}[Hashing a statement by hand]\label{ex:architecture:lock}
Applying the tool's own \lean{strip\_comments} and \lean{statement\_text} functions to
\lean{Epistemic/Tier.lean} gives, for the transitivity theorem, the normalized text
\begin{leancode}
theorem tier_trans {t1 t2 t3 : Tier} (h12 : t1 ≤ t2) (h23 : t2 ≤ t3) : t1 ≤ t3
\end{leancode}
whose SHA-256 begins \lean{94394facb0c97f29}. Everything after \lean{:=}---the two
\lean{have} lines and \lean{omega}---is absent from the hash. If an agent replaced the
hypothesis \lean{(h23 : t2 ≤ t3)} by \lean{(h23 : t1 ≤ t3)}, making the theorem trivially
provable, the normalized text would change and so would the hash; the check would print
\texttt{CHANGED} and exit 1. If instead the agent edited \lean{Tier.toNat} so that
\lean{Tier.A} mapped to $3$ and \lean{Tier.B} to $4$, no theorem statement would change,
but the definition's hash (\texttt{d839f1061e526f25} on the current text) would, and the
check would fail for the same reason: every theorem about $\le$ on tiers now means something
else. This is why definitions are hashed whole and theorems only up to \lean{:=}.
\end{example}

\subsection{Producer $\ne$ verifier}\label{sec:architecture:producer}

The last gate is a rule about people and programs rather than a script: ``No agent or model
grades its own output; the gate runs separately'' \source{STREAM2\_WORKFLOW.md §2,
ll.~25--26}. The measured justification is in the pilot results. Two reports from the
cheapest model tier ``misstated compile status: one claimed 5 clean \lean{sorry}s but left
8 compile errors; one claimed `no compilation errors' with 3 errors. Both caught only because
the orchestrator recompiles every report'' \source{STREAM2\_WORKFLOW.md §6, ll.~153--157};
in the final round a report ``again claimed `0 errors' while one proof failed to parse''
\source{STREAM2\_WORKFLOW.md §6, ll.~192--194}. None of these errors entered the repository,
because the verifier's compile, not the producer's report, is the accept event.

The same principle applies to the project's own tooling. Lesson 0.5 of the log records that
an engine function \lean{sync\_epistemic\_claims} ``injects the same 3 hardcoded claims into
\lean{ledger.jsonl} any time the whole corpus builds with aggregate zero-sorry, regardless of
whether those 3 specific claims are what was actually just proved'', and that a
``prover'' module was a static list with no proving logic \source{LL.md §0.5,
ll.~201--216}. The workflow's rule that ledger entries ``pin a commit hash and build job
count'' is the direct response \source{STREAM2\_WORKFLOW.md §2, ll.~29--31}. A student
reading any auto-generated claim table should ask the question the log asks: which program
wrote this row, and did that program run the kernel?

\begin{pitfallbox}[The gate that no kernel can run]
Gate G2, statement review, is listed in the pipeline as ``the one step no kernel can check:
mis-formalization'' \source{STREAM2\_WORKFLOW.md §4, l.~59}. Two recorded instances show
its two failure modes. Prose drifting past the theorem: a brief said the $\Theta$-shift
preserves $\eta$ ``iff $\Theta$ antisymmetric'' where the Lean proves only ``if''
\source{LL.md §S2.13, ll.~109--113}. And conventions unverified before the statement was
written: every block-matrix sign convention was checked in sympy first, because ``a sign
slip caught this way costs minutes'' and one caught after an agent has spent an hour on an
unprovable statement ``costs far more'' \source{LL.md §S2.12, ll.~101--107}. The lock
freezes the statement; it does not make it right. Only reading does.
\end{pitfallbox}

\subsection{A worked audit: the numbers}\label{sec:architecture:numbers}

The two numbers that this book quotes for the size of the corpus are 326 and 99 audited
theorems. They are the number of \lean{theorem}/\lean{lemma} declarations the audit's
collector finds after stripping comments, every one of which returned an axiom list within
the standard three on a green build. (A naive \lean{grep} for \lean{\textasciicircum
theorem} gives slightly different numbers, because it also counts declarations quoted
inside comments; the author re-ran the audit's own collector on the checkout and obtained
the seven numbers below exactly.) Per library the brief reports
\source{PAPER\_REVISION\_BRIEF §1, ll.~27--35}:
\[
\underbrace{23}_{\text{Validation}}+\underbrace{53}_{\text{Lean5Corpus}}
+\underbrace{44}_{\text{DFT}}+\underbrace{56}_{\text{Foundation}}
+\underbrace{61}_{\text{M24Form.}}+\underbrace{89}_{\text{STFormalization}}=326,
\qquad 326+99=425 .
\]
The reader can reproduce the sum and should. Three cautions attach to the number. It counts
theorems, not physics: the four bridge modules of \cref{sec:architecture:bridges}
contribute 13 of the 56 theorems of the oldest library. It
depends on the build being fresh: the 89 for \lean{StringTheoryFormalization} was reported
as ``89 failing (MISSING)'' once, on a stale \lean{.olean} \source{LL.md §S2.11,
ll.~94--96}. And it is a count \emph{of this plan}: the brief insists that Stream 2's
``18 of 18 planned items'' is ``100\% of that plan, not of string theory''
\source{PAPER\_REVISION\_BRIEF §1, ll.~33--35}. The same brief forbids the phrases ``zero
axioms'', ``100\% mathematical rigor'' and ``formalizes string theory''
\source{PAPER\_REVISION\_BRIEF §0, ll.~15--19}. The correct form, which this book uses
throughout, is: proved in Lean 4 with no axioms beyond the three standard ones.

%=====================================================================================
\section{The tier system, formalized}\label{sec:architecture:epistemic}

The badges \TA, \TL{} and \TC{} used throughout this book are three of five tiers that
the project defines in Lean itself, in \lean{DualScaleM24Formalization/Epistemic/Tier.lean}
and \lean{Ledger.lean}. The two files are short, Mathlib-free, and---unusually for the
Mathlib-free core---not shadows: they state and prove a general theorem about an arbitrary
index type. They are worth reading in full because they make precise a rule that is
otherwise only a convention: \emph{a claim may never be filed at a higher tier than the
claims it rests on.}

\subsection{Tiers as an inductive type}

The tier is an inductive type with five constructors, ordered by a map to $\N$
\source{Tier.lean, ll.~16--38}:
\begin{leancode}
inductive Tier : Type where
  | X : Tier  -- Exploratory / Float
  | C : Tier  -- Conjecture
  | L : Tier  -- Literature
  | B : Tier  -- Exact-Arithmetic Checkable
  | A : Tier  -- Kernel-Verified (Lean 4, zero sorry)
deriving DecidableEq, Repr

def Tier.toNat : Tier → Nat
  | Tier.X => 0 | Tier.C => 1 | Tier.L => 2 | Tier.B => 3 | Tier.A => 4

def Tier.le (t1 t2 : Tier) : Prop := t1.toNat ≤ t2.toNat
\end{leancode}
The module header glosses the five: A is kernel-verified; B is ``deterministic
$\Q/\Z$ check + negative control that fails''; L is literature; C is conjecture; X is
``floats, sampling, LLM output'' \source{Tier.lean, ll.~5--11}. The book collapses B into
symbolic checks and never awards it, and uses X nowhere; the ordering $X<C<L<B<A$ is what
matters. Six theorems establish that $\le$ is a partial order with top element A and bottom
element X, and that nothing but A sits above A \source{Tier.lean, ll.~40--65}. Their
proofs are one-liners: \lean{omega} for the arithmetic ones, \lean{cases t <;> decide} for
the finite case analyses.

\subsection{Ledgers, soundness and transitive support}

A ledger is a function from an index type $\iota$ (the claim identifiers) to rows, each row
carrying a tier and the list of identifiers it cites \source{Ledger.lean, ll.~16--27}:
\begin{leancode}
structure Claim (ι : Type) where
  tier : Tier
  supports : List ι

abbrev Ledger (ι : Type) := ι → Claim ι

def Sound {ι : Type} (L : Ledger ι) : Prop :=
  ∀ a b, b ∈ (L a).supports → (L a).tier ≤ (L b).tier
\end{leancode}
Soundness says: for every row $a$ and every $b$ it cites directly, $\text{tier}(a)\le
\text{tier}(b)$. The interesting content is that the same then holds along chains.
Transitive dependence is an inductive predicate \source{Ledger.lean, ll.~29--43}:
\begin{leancode}
inductive Depends {ι : Type} (L : Ledger ι) : ι → ι → Prop where
  | direct {a b} : b ∈ (L a).supports → Depends L a b
  | step {a b c} : b ∈ (L a).supports → Depends L b c → Depends L a c
\end{leancode}
and the main theorem is proved by induction on it: in a sound ledger, if $a$ depends on $b$
through any chain then $\text{tier}(a)\le\text{tier}(b)$. Two corollaries follow
\source{Ledger.lean, ll.~45--68}.
\begin{itemize}[leftmargin=*,itemsep=1pt]
\item \lean{no\_kernel\_claim\_rests\_on\_weaker}, the one that names the rule: if $a$ is
  filed at A and depends on $b$, then $b$ is at A---because $A\le\text{tier}(b)$ forces
  $\text{tier}(b)=A$ by \lean{eq\_A\_of\_A\_le}.
\item \lean{not\_A\_of\_weak\_support}, its contrapositive and the form one uses in
  practice: one weak link anywhere in the support chain, and the citing claim is not Tier A.
\end{itemize}

Let us do the induction on paper, since it is the whole proof. Suppose $L$ is sound. Case
\lean{direct}: $b\in\text{supports}(a)$, so soundness gives $\text{tier}(a)\le\text{tier}(b)$
at once. Case \lean{step}: $b\in\text{supports}(a)$ and, by the induction hypothesis on
$\text{Depends}\,L\,b\,c$, $\text{tier}(b)\le\text{tier}(c)$; soundness gives
$\text{tier}(a)\le\text{tier}(b)$, and \lean{tier\_trans} chains the two. That is the
entire Lean proof \source{Ledger.lean, ll.~47--52}: the inductive predicate builds the
chain $a=x_0,\ldots,x_k=b$ into the proof principle, so no list ever has to be handled.

\subsection{Positive and negative controls}

The file ends with two concrete ledgers \source{Ledger.lean, ll.~70--91}. The first has
three rows over \lean{Fin 3}: row 0 at A citing row 1, row 1 at A citing nothing, row 2 at C
citing nothing; \lean{witnessSound\_sound} proves it sound by \lean{decide}, which checks
the finitely many pairs $(a,b)$. The second has two rows: row 0 at A citing row 1, row 1 at
C; \lean{witnessUnsound\_not\_sound} proves it is \emph{not} sound---the instance $a=0,b=1$
of the soundness condition would require $A\le C$, which \lean{decide} refutes. This is the
``negative control that fails'' of tier B applied to the tier system itself: a definition of
soundness that accepted the second ledger would be vacuous, and the theorem shows it does
not.

\begin{physicsbox}[What the ledger theorem does and does not say]
The theorem is about an abstract function $\iota\to\text{Claim}\,\iota$. The project also
keeps a concrete file \lean{ledger.jsonl} with, on the current checkout, 38 rows (34 filed at
A, 3 at L, 1 at B), each naming a Lean source file and a list of dependencies, together with
a Markdown rendering \lean{LEDGER.md} whose header quotes the transitivity theorem. Two
things a student must keep apart. First, no Lean theorem says that \lean{ledger.jsonl} is a
\lean{Sound} ledger: the JSON is not a term of type \lean{Ledger ι}, and the ``dependencies''
column is bookkeeping, not imports---eleven of the Stream 1 rows list the ledger theorem
\lean{K3T2-A-0001} itself as their dependency, while the source file behind one of them,
\lean{EffectiveMetric.lean}, has no \lean{import} line at all. Second, the \emph{statement
text} of a row can promise more than the Lean it points to. The row \lean{OBS-VAL-0001}
is filed at A with the words ``Certified Falsifiable Observables \ldots\ LiteBIRD $r =
0.00396$''; the module it names proves that $1/252$ lies in the window $[0.003,0.005]$ and
similar inequalities between literals, and its own header says it is ``Tier A for its
arithmetic only'' with the physical relations Tier C \source{Observables.lean, ll.~9--24}.
The ledger theorem guarantees nothing about how a row is worded. The tier badge in this
book therefore attaches to the Lean statement, never to a ledger row.
\end{physicsbox}

\begin{remark}[What a faithful ledger would need]
To make the theorem bite on the real corpus one would generate the \lean{Ledger} term from
the declaration graph (\cref{ch:atlas}): identifiers are declaration names, \lean{supports}
is the list of project declarations a proof term uses, and \lean{tier} is read from a
docstring tag. Soundness would then be decidable by \lean{decide} exactly as for
\lean{witnessSound}, and a Tier A theorem citing a Tier L definition would fail to compile.
Nothing in the present files does this; it is one of the improvements the log lists for the
next run (``paper claims come from the Lean source'') \source{LL.md §S2-I, ll.~139--141}.
\end{remark}

%=====================================================================================
\section{What the machine has checked}\label{sec:architecture:lean}

This chapter's Lean content is the epistemic module itself; the theorems of the other
libraries are the subject of their own chapters. Every statement below was read from the
source file and audited by the author of this chapter as in \cref{ex:architecture:audit}.

\begin{leanbox}[The tier order]
\begin{leancode}
theorem tier_refl (t : Tier) : t ≤ t := by ...
theorem tier_trans {t1 t2 t3 : Tier} (h12 : t1 ≤ t2) (h23 : t2 ≤ t3) :
    t1 ≤ t3 := by ...
theorem tier_antisymm {t1 t2 : Tier} (h12 : t1 ≤ t2) (h21 : t2 ≤ t1) :
    t1 = t2 := by ...
theorem tier_A_is_maximal (t : Tier) : t ≤ Tier.A := by ...
theorem tier_X_is_minimal (t : Tier) : Tier.X ≤ t := by ...
theorem eq_A_of_A_le {s : Tier} (h : Tier.A ≤ s) : s = Tier.A := by ...
\end{leancode}
\textbf{Reading.} The relation $t_1\le t_2 :\Leftrightarrow \text{toNat}(t_1)\le
\text{toNat}(t_2)$ on the five-element type is reflexive, transitive and antisymmetric;
A is the greatest element, X the least, and the only tier $\ge$ A is A.
\textbf{Proof idea.} Unfold to $\N$ and call \lean{omega}; for antisymmetry, derive equality
of the images and split into the $25$ constructor cases, each closed by \lean{rfl} or
\lean{contradiction} \source{Tier.lean, ll.~40--65}. \textbf{Not proved.} Nothing about
what the tiers mean; the type carries no semantics beyond the order. \TA
\end{leanbox}

\begin{leanbox}[Transitive soundness of a ledger]
\begin{leancode}
theorem tier_le_of_depends {ι : Type} {L : Ledger ι} (hL : Sound L) :
    ∀ {a b : ι}, Depends L a b → (L a).tier ≤ (L b).tier := by ...
theorem no_kernel_claim_rests_on_weaker {ι : Type} {L : Ledger ι} (hL : Sound L)
    {a b : ι} (ha : (L a).tier = Tier.A) (hab : Depends L a b) :
    (L b).tier = Tier.A := by ...
theorem not_A_of_weak_support {ι : Type} {L : Ledger ι} (hL : Sound L)
    {a b : ι} (hab : Depends L a b) (hb : (L b).tier ≠ Tier.A) :
    (L a).tier ≠ Tier.A := ...
theorem witnessSound_sound : Sound witnessSound := by ...
theorem witnessUnsound_not_sound : ¬ Sound witnessUnsound := by ...
\end{leancode}
\textbf{Reading.} For every index type and every ledger satisfying the one-step condition
\lean{Sound}: tiers are monotone along every dependency chain; a Tier A row's entire
transitive support is Tier A; a row with any non-A support is not A. The two witnesses are
a three-row sound ledger and a two-row unsound one. \textbf{Proof idea.} Induction on the
inductive predicate \lean{Depends}, using soundness at each step and \lean{tier\_trans};
the corollary rewrites with \lean{ha} and applies \lean{eq\_A\_of\_A\_le}; the witnesses
are finite and closed by \lean{decide} \source{Ledger.lean, ll.~45--91}. \textbf{Not
proved.} That any actual file (\lean{ledger.jsonl}, the declaration graph) is a
\lean{Sound} ledger; that a row's wording matches its source. \TA
\end{leanbox}

\begin{center}\small
\begin{tabular}{@{}p{7.4cm}p{1.2cm}p{5.6cm}@{}}\toprule
Claim & Tier & Where \\ \midrule
Tier order: partial order, A maximal, X minimal & \TA & \lean{Tier.lean} \\
Transitive tier monotonicity in a sound ledger; A rests only on A & \TA & \lean{Ledger.lean} \\
Positive and negative control ledgers & \TA & \lean{Ledger.lean} \\
Axiom footprint of the five theorems above: $\subseteq\{\lean{propext},\lean{Quot.sound}\}$ & \TA & \cref{ex:architecture:audit} \\
$(\eta\HH)^2=1$ for all invertible real $G$, all $B$, all $d$ & \TA & \lean{DFT/GeneralizedMetric.lean} \\
Burnside guard $\sum\dim^2=|M_{24}|$ on the corrected table & \TA & \lean{MathieuM24.lean} \\
Burnside's identity itself; $H^2(K3,\Z)\cong E_8(-1)^2\oplus U^3$ & \TL & standard; Huybrechts, via the brief \\
\lean{ledger.jsonl} is a sound ledger; any row's physical wording & --- & not formalized; not a Lean claim \\
The seven-library split and gate sequence are the right architecture & \TC & this programme's design choice \\
\bottomrule\end{tabular}
\end{center}

\begin{summarybox}
\begin{itemize}[leftmargin=*,itemsep=1pt]
\item One Lake project, Lean and Mathlib pinned to \lean{v4.33.1}, seven libraries: five
  Mathlib-free default targets (\lean{StringTheoryFoundation},
  \lean{DualScaleM24Formalization}, \lean{DoubleFieldTheory}, \lean{DualScaleValidation},
  \lean{Lean5Corpus}) and two Mathlib-based libraries built by name
  (\lean{StringTheoryFormalization}, \lean{DualScaleStream2}). Imports run one way:
  Stream 2 imports Stream 1, never the reverse.
\item Audited on 2026-09-17: $326$ theorems in Stream 1 ($23+53+44+56+61+89$) and $99$ in
  Stream 2, every one depending only on \lean{propext}, \lean{Classical.choice},
  \lean{Quot.sound}.
\item Depth differs by library. The Mathlib-free core mostly proves arithmetic shadows
  (integer and rational instances, fixed small matrices); the Mathlib-based libraries prove
  the general statements (real $d\times d$ matrices, every $d$). Quote the statement as
  written, and name its library.
\item Four bridge modules named after external formalizations import nothing from them and
  prove thirteen integer identities; each says so in a scope note.
\item Gates: build, \lean{sorry} count, statement review, axiom audit
  (\lean{\#print axioms}, catching \lean{sorryAx} and \lean{Lean.ofReduceBool}), statement
  lock (SHA-256 of the statement up to \lean{:=}, of definitions whole), producer $\ne$
  verifier, and ledger rows pinned to a commit. Each exists because of a recorded failure.
\item The tier system is itself formalized: in a \lean{Sound} ledger a Tier A claim rests
  only on Tier A claims (\lean{no\_kernel\_claim\_rests\_on\_weaker}), a theorem about an
  abstract ledger type, not about any file.
\end{itemize}
\end{summarybox}

%=====================================================================================
\section*{Exercises}\addcontentsline{toc}{section}{Exercises}

\begin{exercise}[$\star$]
Reproduce the audit totals: add the six Stream 1 counts of
\cref{tab:architecture:libraries} and check that they give 326, then add Stream 2. Which
library contributes the most theorems per module, and which the fewest? (Use the module
counts in the table.)
\end{exercise}

\begin{exercise}[$\star$]
From \cref{fig:architecture:imports}, list every library that can be compiled without
Mathlib installed. Explain why a compile error in \lean{DualScaleStream2} cannot break
\lean{Lean5Corpus}, and why an error in \lean{StringTheoryFoundation} can break four
other libraries.
\end{exercise}

\begin{exercise}[$\star\star$]
Explain precisely why \lean{gen\_metric\_duality} (hypothesis $g=1$, integer entries) does
not imply \lean{etaR\_genMetric\_sq} even for $d=1$. Then write, in Mathlib notation, the
$d=1$ real statement: for $G=(g)$ with $g\neq0$ and $B=(0)$, the matrix $\eta\HH$ squares to
the identity. Check it by hand: $\HH=\diag(g,g^{-1})$, $\eta\HH=\begin{psmallmatrix}0&g^{-1}\\
g&0\end{psmallmatrix}$, and $(\eta\HH)^2=\mathbf 1$.
\end{exercise}

\begin{exercise}[$\star\star$]
For each edit, say whether \lean{statement\_lock.py --check} would report \texttt{CHANGED},
and why: (a) replacing \lean{:= by sorry} with \lean{:= by omega} in \lean{tier\_trans};
(b) adding the hypothesis \lean{(h : t1 = t2)} to \lean{tier\_antisymm}; (c) rewriting the
docstring of \lean{Sound}; (d) changing \lean{Tier.toNat} so that \lean{Tier.L} maps to $3$
and \lean{Tier.B} to $2$; (e) reformatting \lean{Depends} onto one line.
\end{exercise}

\begin{exercise}[$\star\star$]
Which of the following \lean{\#print axioms} outputs pass the audit of
\cref{sec:architecture:axioms}, and what does each failing one indicate?
\begin{enumerate}[label=(\alph*),leftmargin=*,itemsep=0pt]
\item \lean{[propext]};
\item \lean{[propext, Classical.choice, Quot.sound]};
\item \lean{[propext, sorryAx]};
\item \lean{[Lean.ofReduceBool]};
\item ``does not depend on any axioms''.
\end{enumerate}
\end{exercise}

\begin{exercise}[$\star\star$ Lean]
In a file importing \lean{DualScaleM24Formalization.Epistemic.Tier}, state and prove
\begin{leancode}
theorem B_le_A : Tier.B ≤ Tier.A := by sorry
theorem not_A_le_C : ¬ (Tier.A ≤ Tier.C) := by sorry
\end{leancode}
(Hint: both are closed by \lean{decide}, since \lean{≤} on \lean{Tier} was given a
\lean{Decidable} instance \source{Tier.lean, ll.~37--38}.)
\end{exercise}

\begin{exercise}[$\star\star\star$ Lean]
Define a four-row ledger over \lean{Fin 4} in which row 0 (tier A) cites row 1, row 1
(tier A) cites row 2, row 2 (tier L) cites row 3 and row 3 is tier C. Prove it is not
\lean{Sound}, and identify the offending pair. Then modify the tiers minimally so that it is
sound, prove soundness by \lean{decide}, and use \lean{tier\_le\_of\_depends} to conclude
that the tier of row 0 is at most that of row 3.
\end{exercise}

\begin{exercise}[$\star\star\star$]
The theorem \lean{energy\_dissipation\_monotonic (e d : Nat) : e - d ≤ e} uses truncated
subtraction, so $e-d=0$ whenever $d>e$. State the energy inequality
$\frac{\dd}{\dd t}\|u\|_{L^2}^2=-2\nu\|\nabla u\|_{L^2}^2\le0$ that the module header
describes \source{NavierStokesBridge.lean, ll.~17--37}, and list the objects a faithful Lean
statement would need (a function space, a time derivative, an integral, a hypothesis that
$u$ solves the equation). Which of them does Mathlib provide, and why can none of them
appear in a Mathlib-free library?
\end{exercise}

\begin{exercise}[$\star\star$]
Eleven rows of \lean{ledger.jsonl} list \lean{K3T2-A-0001} (the ledger theorem itself) as
their sole dependency. Is this consistent with \lean{Sound}? Is it a mathematical
dependency? Propose a rule for what ``dependency'' should mean in a machine-generated
ledger, and say how the declaration graph of \cref{ch:atlas} would supply it.
\end{exercise}

%=====================================================================================
\section*{Further reading}\addcontentsline{toc}{section}{Further reading}
\begin{itemize}[leftmargin=*,itemsep=2pt]
\item The build description itself: \lean{lakefile.lean} \source{ll.~1--53}, in particular
  the comments on the Mathlib pin and on why the two Mathlib-based libraries are not default
  targets.
\item The workflow that produced Stream 2: hard constraints \source{STREAM2\_WORKFLOW.md §2,
  ll.~20--38}, the per-package gates \source{§4, ll.~52--65}, and the final gate numbers
  \source{§6, ll.~185--189}. \Cref{ch:pipeline} discusses the measured tier split.
\item The verified-facts brief that fixes the vocabulary and the audit numbers used in this
  book \source{PAPER\_REVISION\_BRIEF §0--§1, ll.~8--35} and the corrections that must
  appear in the papers \source{§2, ll.~37--62}.
\item The lessons behind each gate: recovery and build-fix \source{LL.md §0, ll.~145--259},
  the axiom audit \source{§S2.5, ll.~56--61}, guard theorems \source{§S2.6, ll.~63--67},
  the statement lock \source{§S2.10, ll.~87--92}, failing loudly \source{§S2.11,
  ll.~94--99}, and the planned single gate script \source{§S2-I, ll.~116--143}.
\item The tools \lean{tools/axiom\_audit.py} \source{ll.~1--102} and
  \lean{tools/statement\_lock.py} \source{ll.~1--137}, both under 150 lines, and the two
  epistemic modules \lean{Epistemic/Tier.lean} \source{ll.~1--67} and
  \lean{Epistemic/Ledger.lean} \source{ll.~1--93}; a student can read all four in an evening.
\end{itemize}
